(a) If $D=\sum n_PP$ is effective, the sheaf $\mathcal O_D$ has, at each $P$, a filtration with $n_P$ quotients isomorphic to $k(P)$. Thus $\gamma(\mathcal O_D)=\psi(D)$. For any divisor $D$ and closed point $P$, (6.18), tensored with $\mathcal L(D)$, gives $0\to\mathcal L(D-P)\to\mathcal L(D)\to k(P)\to0$. Adding and subtracting points therefore gives
\[\psi(D)=\gamma(\mathcal L(D))-\gamma(\mathcal O_X).\]
Since $\mathcal L(D)$ depends only on the linear equivalence class, $\psi$ descends to $\operatorname{Cl}X$.

(b) Embed $X$ in its nonsingular projective completion by (I, 6.10), and extend $\mathcal F$ coherently by \exref{II.5.15}. Applying (5.17) after twisting and restricting back to $X$ gives a surjection $\mathcal E_0\to\mathcal F$ with $\mathcal E_0$ locally free. Its kernel $\mathcal E_1$ is coherent and torsion-free, hence locally free since the local rings at closed points are DVRs.

For two such resolutions, the fiber product $\mathcal B=\mathcal E_0\times_{\mathcal F}\mathcal E'_0$ fits into $0\to\mathcal E_1\to\mathcal B\to\mathcal E'_0\to0$ and $0\to\mathcal E'_1\to\mathcal B\to\mathcal E_0\to0$. These are locally split sequences of locally free sheaves. By \exref{II.5.16}(d),
\[\det\mathcal E_1\otimes\det\mathcal E'_0\cong\det\mathcal B\cong\det\mathcal E'_1\otimes\det\mathcal E_0.\]
Thus $\det\mathcal F=\det\mathcal E_0\otimes(\det\mathcal E_1)^{-1}$ is independent of the resolution.

Given $0\to\mathcal F'\to\mathcal F\to\mathcal F''\to0$, choose a locally free surjection $\mathcal E_0\to\mathcal F$. Put $\mathcal K=\ker(\mathcal E_0\to\mathcal F)$ and $\mathcal H=\ker(\mathcal E_0\to\mathcal F'')$, both locally free. The sequences resolving $\mathcal F,\mathcal F',\mathcal F''$ have pairs $(\mathcal E_0,\mathcal K)$, $(\mathcal H,\mathcal K)$, and $(\mathcal E_0,\mathcal H)$, respectively. Their determinant formulas give $\det\mathcal F\cong\det\mathcal F'\otimes\det\mathcal F''$. Hence determinant induces a homomorphism on $K(X)$. Finally, $0\to\mathcal O_X(-P)\to\mathcal O_X\to k(P)\to0$ gives $\det\gamma(k(P))=\mathcal O_X(P)$, so $\det\psi(D)=\mathcal L(D)$.

(c) Choose a basis of $\mathcal F$ at the generic point. It extends to sections on some nonempty open set. Choose an effective divisor supported on the complement. On a finite affine cover trivializing this divisor, (5.3) extends the sections after multiplying by powers of its local equations. Their differences on overlaps vanish on the original open set, so a further common power kills these differences by (5.3). Thus, for a sufficiently large effective divisor $E$, they give a map $\mathcal O_X(-E)^{\oplus r}\to\mathcal F$ which is an isomorphism at the generic point. Its kernel is a torsion subsheaf of a locally free sheaf, hence zero, and its cokernel $\mathcal T$ is torsion. Thus the required divisor is $D=-E$. A composition series at each point gives $\gamma(\mathcal T)=\sum_P\operatorname{length}(\mathcal T_P)\gamma(k(P))$. Together with (a), this yields
\[\gamma(\mathcal F)-r\gamma(\mathcal O_X)=\psi\left(-rE+\sum_P\operatorname{length}(\mathcal T_P)P\right).\]

(d) The map $(\det,\operatorname{rank}):K(X)\to\operatorname{Pic}X\oplus\mathbb Z$ is an isomorphism. Indeed, its inverse sends $(\mathcal L(D),r)$ to $\psi(D)+r\gamma(\mathcal O_X)$. The two composites are identities by (b),(c), since $\operatorname{rank}\psi(D)=0$ and $\det\mathcal O_X=\mathcal O_X$.
